% -*- coding: utf-8 -*-
% !TEX program = xelatex
\documentclass[plain,noanswer,medmath]{../../template/jnuexam}
\usepackage{../../template/kysx}

\begin{document}


\examtitle{2017}{1}


\makepart{选择题}{1～8小题，每小题4分，共32分}


\begin{problem}
若函数$f(x) = \begin{cases}
  \frac{1 - \cos \sqrt{x}}{ax}, & x > 0 \\
                             b, & x \le 0
\end{cases}$在$x = 0$处连续，则\pickout{}
\begin{abcd}
\item $ab = \frac{1}{2}$．
\item $ab = - \frac{1}{2}$．
\item $ab = 0$．
\item $ab = 2$．
\end{abcd}
\end{problem}


\begin{problem}
若函数$f(x)$可导，且$f(x)f'(x) > 0$，则\pickout{}
\begin{abcd}
\item $f(1) > f(-1)$．
\item $f(1) < f(-1)$．
\item $\left| f(1) \right| > \left| f(-1) \right|$．
\item $\left| f(1) \right| < \left| f(-1) \right|$．
\end{abcd}
\end{problem}


\begin{problem}
函数$f(x,y,z) = x^2 y + z^2$在点$(1,2,0)$处沿向量$u = (1,2,2)$的方向导数为\pickout{}
\begin{abcd}
\item $12$．
\item $6$．
\item $4$．
\item $2$．
\end{abcd}
\end{problem}


\begin{problem}
甲、乙两人赛跑，计时开始时，甲在乙前方$10$（单位：m）处，
图中实线表示甲的速度曲线$v = v_1(t)$（单位：m/s），虚线表示乙的速度$v = v_2(t)$，
三块阴影部分面积的数值依次为$10$、$20$、$3$，计时开始后乙追上甲的时刻记为$t_0$（单位：s），
则\pickout{}
\par\centerline{\includegraphics[width=0.5\linewidth]{2017A}}
\begin{abcd}
\item $t_0 = 10$．
\item $15 < t_0 < 20$．
\item $t_0 = 25$．
\item $t_0 > 25$．
\end{abcd}
\end{problem}


\begin{problem}
设$\alpha$是$n$维单位列向量，$E$为$n$阶单位矩阵，则\pickout{}
\begin{abcd}
\item $E - \alpha \alpha^T$不可逆．
\item $E + \alpha \alpha^T$不可逆．
\item $E + 2\alpha \alpha^T$不可逆．
\item $E - 2\alpha \alpha^T$不可逆．
\end{abcd}
\end{problem}


\begin{problem}
设矩阵$A = \begin{pmatrix}
  2 & 0 & 0 \\
  0 & 2 & 1 \\
  0 & 0 & 1
\end{pmatrix}$，$B = \begin{pmatrix}
  2 & 1 & 0 \\
  0 & 2 & 0 \\
  0 & 0 & 1
\end{pmatrix}$，$C = \begin{pmatrix}
  1 & 0 & 0 \\
  0 & 2 & 0 \\
  0 & 0 & 2
\end{pmatrix}$，则\pickout{}
\begin{abcd}
\item $A$与$C$相似，$B$与$C$相似．
\item $A$与$C$相似，$B$与$C$不相似．
\item $A$与$C$不相似，$B$与$C$相似．
\item $A$与$C$不相似，$B$与$C$不相似．
\end{abcd}
\end{problem}


\begin{problem}
设$A,B$为随机概率，若$0 < P(A) < 1$，$0 < P(B) < 1$，
则$P(A\left| B \right.) > P(A\left| \bar B \right.)$的充分必要条件是\pickout{}
\begin{abcd}
\item $P(B\left| A \right.) > P(B\left| \bar A \right.)$．
\item $P(B\left| A \right.) < P(B\left| \bar A \right.)$．
\item $P(\bar B\left| A \right.) > P(B\left| \bar A \right.)$．
\item $P(\bar B\left| A \right.) < P(B\left| \bar A \right.)$．
\end{abcd}
\end{problem}


\begin{problem}
设$X_1,X_2 \cdots X_n(n \ge 2)$为来自总体$N(\mu ,1)$的简单随即样本，
记$\overline{X} = \frac{1}{n}\sum_{i = 1}^n X_i$，则下列结论中不正确的是\pickout{}
\begin{abcd}
\item $\sum_{i = 1}^n (X_i - \mu)^2$服从$\chi^2$分布．
\item $2(X_n - X_1)^2$服从$\chi^2$分布．
\item $\sum_{i = 1}^n (X_i - \overline{X})^2$服从$\chi^2$分布．
\item $n(\overline{X} - \mu)^2$服从$\chi^2$分布．
\end{abcd}
\end{problem}


\makepart{填空题}{9～14小题，每小题4分，共24分}


\begin{problem}
已知函数$f(x) = \frac{1}{1 + x^2}$，则$f^{(3)}(0) =$ \fillin{}．
\end{problem}


\begin{problem}
微分方程$y'' + 2y' + 3y = 0$的通解为$y =$ \fillin{}．
\end{problem}


\begin{problem}
若曲线积分$\int_L \frac{x\dx - ay\dy}{x^2 + y^2 - 1}$
在区域$D = \{(x,y) \left| x \right.^2 + y^2 < 1\}$内与路径无关，则$a =$ \fillin{}．
\end{problem}


\begin{problem}
幂级数$\sum_{n = 1}^{\infty}(- 1)^{n - 1} nx^{n - 1}$在区间$(- 1,1)$内的和函数$S(x) =$ \fillin{}．
\end{problem}


\begin{problem}
设矩阵$A = \begin{pmatrix}
  1 & 0 & 1 \\
  1 & 1 & 2 \\
  0 & 1 & 1
\end{pmatrix}$，$\alpha_1$, $\alpha_2$, $\alpha_3$为线性无关的$3$维列向量组，
则向量组$A\alpha_1$, $A\alpha_2$, $A\alpha_3$的秩为 \fillin{}．
\end{problem}


\begin{problem}
设随机变量$X$的分布函数为$F(x) = 0.5\Phi(x) + 0.5\Phi \left(\frac{x - 4}{2} \right)$，
其中$\Phi (x)$为标准正态分布函数，则$EX =$ \fillin{}．
\end{problem}


\makepart{解答题}{15～23小题，共94分}


\begin{problem}[points=10]
设函数$f(u,v)$具有$2$阶连续偏导数，$y = f(\e^x,\cos x)$，求$\frac{\dy}{\dx}\big|_{x = 0}$，$\frac{\d^2y}{\dx^2}\big|_{x = 0}$．
\end{problem}


\begin{problem}[points=10]
求$\lim_{n \to \infty} \sum_{k = 1}^n \frac{k}{n^2} \ln \left(1 + \frac{k}{n} \right)$．
\end{problem}


\begin{problem}[points=10]
已知函数$y(x)$由方程$x^3 + y^3 - 3x + 3y - 2 = 0$确定，求$y(x)$的极值．
\end{problem}


\begin{problem}[points=11]
设函数$f(x)$在区间$[0,1]$上具有二阶导数，且$f(1) > 0$，$\lim_{x \to 0^+} \frac{f(x)}{x} < 0$．证明：
\begin{enumerate}
\item 方程$f(x) = 0$在区间$(0,1)$内至少存在一个实根．
\item 方程$f(x)f''(x) + \left(f'(x) \right)^2 = 0$在区间$(0,1)$内至少存在两个不同实根．
\end{enumerate}
\end{problem}


\begin{problem}[points=10]
设薄片型物体$S$是圆锥面$z = \sqrt{x^2 + y^2}$被柱面$z^2 = 2x$割下的有限部分，
其上任一点的密度为$\mu = 9\sqrt{x^2 + y^2 + z^2}$，记圆锥面与柱面的交线为$C$．\par
\begin{enumerate*}
\item 求$C$在$xOy$面上的投影曲线的方程；
\item 求$S$的质量$M$．
\end{enumerate*}
\end{problem}


\begin{problem}[points=11]
设$3$阶矩阵$A= (\alpha_1,\alpha_2,\alpha_3)$有$3$个不同的特征值，
且$\alpha_3=\alpha_1\text{+}2\alpha_2$．\par
\begin{enumerate*}
\item 证明：$r(A) = 2$；
\item 若$\beta =\alpha_1\text{+}\alpha_2\text{+}\alpha_3$，求方程组$Ax = \beta$的通解．
\end{enumerate*}
\end{problem}


\begin{problem}[points=11]
设二次型$f(x_1,x_2,x_3) = 2x_1^2 - x_2^2 + ax_3^2 + 2x_1x_2 - 8x_1x_3 + 2x_2x_3$
在正交变换$x = Qy$下的标准型为$\lambda_1 y_1^2 + \lambda_2 y_2^2$，求$a$的值及一个正交矩阵$Q$．
\end{problem}


\begin{problem}[points=11]
设随机变量为$X$，$Y$相互独立，且$X$的概率分布为$P(X = 0) = P(X = 2) = \frac{1}{2}$，
Y的概率密度为$f(y) = \begin{cases}
  2y, & 0 < y < 1, \\
   0, & \text{其他}.
\end{cases}$\par
\begin{enumerate*}
\item 求$P(Y \le EY)$；
\item 求$Z = X + Y$的概率密度．
\end{enumerate*}
\end{problem}


\begin{problem}[points=11]
某工程师为了解一台天平的精度，用该天平对一物体的质量做$n$次测量，该物体的质量$\mu$是已知的，
设$n$次测量结果$X_1,X_2, \cdots ,X_n$相互独立且均服从正态分布$N(\mu ,\sigma^2)$．
该工程师记录的是$n$次测量的绝对误差$Z_i = |X_i-\mu|$（$i = 1,2, \cdots,n$），
利用$Z_1,Z_2, \cdots ,Z_n$估计$\sigma$．
\begin{enumerate}
\item 求$Z_i$的概率密度；
\item 利用一阶矩求$\sigma$的矩估计量；
\item 求$\sigma$的最大似然估计量．
\end{enumerate}
\end{problem}


\end{document}
